\IfFileExists{stacks-project.cls}{%
\documentclass{stacks-project}
}{%
\documentclass{amsart}
}

% For dealing with references we use the comment environment
\usepackage{verbatim}
\newenvironment{reference}{\comment}{\endcomment}
%\newenvironment{reference}{}{}
\newenvironment{slogan}{\comment}{\endcomment}
\newenvironment{history}{\comment}{\endcomment}

% For commutative diagrams we use Xy-pic
\usepackage[all]{xy}

% We use 2cell for 2-commutative diagrams.
\xyoption{2cell}
\UseAllTwocells

% We use multicol for the list of chapters between chapters
\usepackage{multicol}

% This is generally recommended for better output
\usepackage{lmodern}
\usepackage[T1]{fontenc}

% For cross-file-references
\usepackage{xr-hyper}

% Package for hypertext links:
\usepackage{hyperref}

% For any local file, say "hello.tex" you want to link to please
% use \externaldocument[hello-]{hello}
\externaldocument[introduction-]{introduction}
\externaldocument[conventions-]{conventions}
\externaldocument[sets-]{sets}
\externaldocument[categories-]{categories}
\externaldocument[topology-]{topology}
\externaldocument[sheaves-]{sheaves}
\externaldocument[sites-]{sites}
\externaldocument[stacks-]{stacks}
\externaldocument[fields-]{fields}
\externaldocument[algebra-]{algebra}
\externaldocument[brauer-]{brauer}
\externaldocument[homology-]{homology}
\externaldocument[derived-]{derived}
\externaldocument[simplicial-]{simplicial}
\externaldocument[more-algebra-]{more-algebra}
\externaldocument[smoothing-]{smoothing}
\externaldocument[modules-]{modules}
\externaldocument[sites-modules-]{sites-modules}
\externaldocument[injectives-]{injectives}
\externaldocument[cohomology-]{cohomology}
\externaldocument[sites-cohomology-]{sites-cohomology}
\externaldocument[dga-]{dga}
\externaldocument[dpa-]{dpa}
\externaldocument[sdga-]{sdga}
\externaldocument[hypercovering-]{hypercovering}
\externaldocument[schemes-]{schemes}
\externaldocument[constructions-]{constructions}
\externaldocument[properties-]{properties}
\externaldocument[morphisms-]{morphisms}
\externaldocument[coherent-]{coherent}
\externaldocument[divisors-]{divisors}
\externaldocument[limits-]{limits}
\externaldocument[varieties-]{varieties}
\externaldocument[topologies-]{topologies}
\externaldocument[descent-]{descent}
\externaldocument[perfect-]{perfect}
\externaldocument[more-morphisms-]{more-morphisms}
\externaldocument[flat-]{flat}
\externaldocument[groupoids-]{groupoids}
\externaldocument[more-groupoids-]{more-groupoids}
\externaldocument[etale-]{etale}
\externaldocument[chow-]{chow}
\externaldocument[intersection-]{intersection}
\externaldocument[pic-]{pic}
\externaldocument[weil-]{weil}
\externaldocument[adequate-]{adequate}
\externaldocument[dualizing-]{dualizing}
\externaldocument[duality-]{duality}
\externaldocument[discriminant-]{discriminant}
\externaldocument[derham-]{derham}
\externaldocument[local-cohomology-]{local-cohomology}
\externaldocument[algebraization-]{algebraization}
\externaldocument[curves-]{curves}
\externaldocument[resolve-]{resolve}
\externaldocument[models-]{models}
\externaldocument[functors-]{functors}
\externaldocument[equiv-]{equiv}
\externaldocument[pione-]{pione}
\externaldocument[etale-cohomology-]{etale-cohomology}
\externaldocument[proetale-]{proetale}
\externaldocument[relative-cycles-]{relative-cycles}
\externaldocument[more-etale-]{more-etale}
\externaldocument[trace-]{trace}
\externaldocument[crystalline-]{crystalline}
\externaldocument[spaces-]{spaces}
\externaldocument[spaces-properties-]{spaces-properties}
\externaldocument[spaces-morphisms-]{spaces-morphisms}
\externaldocument[decent-spaces-]{decent-spaces}
\externaldocument[spaces-cohomology-]{spaces-cohomology}
\externaldocument[spaces-limits-]{spaces-limits}
\externaldocument[spaces-divisors-]{spaces-divisors}
\externaldocument[spaces-over-fields-]{spaces-over-fields}
\externaldocument[spaces-topologies-]{spaces-topologies}
\externaldocument[spaces-descent-]{spaces-descent}
\externaldocument[spaces-perfect-]{spaces-perfect}
\externaldocument[spaces-more-morphisms-]{spaces-more-morphisms}
\externaldocument[spaces-flat-]{spaces-flat}
\externaldocument[spaces-groupoids-]{spaces-groupoids}
\externaldocument[spaces-more-groupoids-]{spaces-more-groupoids}
\externaldocument[bootstrap-]{bootstrap}
\externaldocument[spaces-pushouts-]{spaces-pushouts}
\externaldocument[spaces-chow-]{spaces-chow}
\externaldocument[groupoids-quotients-]{groupoids-quotients}
\externaldocument[spaces-more-cohomology-]{spaces-more-cohomology}
\externaldocument[spaces-simplicial-]{spaces-simplicial}
\externaldocument[spaces-duality-]{spaces-duality}
\externaldocument[formal-spaces-]{formal-spaces}
\externaldocument[restricted-]{restricted}
\externaldocument[spaces-resolve-]{spaces-resolve}
\externaldocument[formal-defos-]{formal-defos}
\externaldocument[defos-]{defos}
\externaldocument[cotangent-]{cotangent}
\externaldocument[examples-defos-]{examples-defos}
\externaldocument[algebraic-]{algebraic}
\externaldocument[examples-stacks-]{examples-stacks}
\externaldocument[stacks-sheaves-]{stacks-sheaves}
\externaldocument[criteria-]{criteria}
\externaldocument[artin-]{artin}
\externaldocument[quot-]{quot}
\externaldocument[stacks-properties-]{stacks-properties}
\externaldocument[stacks-morphisms-]{stacks-morphisms}
\externaldocument[stacks-limits-]{stacks-limits}
\externaldocument[stacks-cohomology-]{stacks-cohomology}
\externaldocument[stacks-perfect-]{stacks-perfect}
\externaldocument[stacks-introduction-]{stacks-introduction}
\externaldocument[stacks-more-morphisms-]{stacks-more-morphisms}
\externaldocument[stacks-geometry-]{stacks-geometry}
\externaldocument[moduli-]{moduli}
\externaldocument[moduli-curves-]{moduli-curves}
\externaldocument[examples-]{examples}
\externaldocument[exercises-]{exercises}
\externaldocument[guide-]{guide}
\externaldocument[desirables-]{desirables}
\externaldocument[coding-]{coding}
\externaldocument[obsolete-]{obsolete}
\externaldocument[fdl-]{fdl}
\externaldocument[index-]{index}

% Theorem environments.
%
\theoremstyle{plain}
\newtheorem{theorem}[subsection]{Theorem}
\newtheorem{proposition}[subsection]{Proposition}
\newtheorem{lemma}[subsection]{Lemma}

\theoremstyle{definition}
\newtheorem{definition}[subsection]{Definition}
\newtheorem{example}[subsection]{Example}
\newtheorem{exercise}[subsection]{Exercise}
\newtheorem{situation}[subsection]{Situation}

\theoremstyle{remark}
\newtheorem{remark}[subsection]{Remark}
\newtheorem{remarks}[subsection]{Remarks}

\numberwithin{equation}{subsection}

% Macros
%
\def\lim{\mathop{\mathrm{lim}}\nolimits}
\def\colim{\mathop{\mathrm{colim}}\nolimits}
\def\Spec{\mathop{\mathrm{Spec}}}
\def\Hom{\mathop{\mathrm{Hom}}\nolimits}
\def\Ext{\mathop{\mathrm{Ext}}\nolimits}
\def\SheafHom{\mathop{\mathcal{H}\!\mathit{om}}\nolimits}
\def\SheafExt{\mathop{\mathcal{E}\!\mathit{xt}}\nolimits}
\def\Sch{\mathit{Sch}}
\def\Mor{\mathop{\mathrm{Mor}}\nolimits}
\def\Ob{\mathop{\mathrm{Ob}}\nolimits}
\def\Sh{\mathop{\mathit{Sh}}\nolimits}
\def\NL{\mathop{N\!L}\nolimits}
\def\CH{\mathop{\mathrm{CH}}\nolimits}
\def\proetale{{pro\text{-}\acute{e}tale}}
\def\etale{{\acute{e}tale}}
\def\QCoh{\mathit{QCoh}}
\def\Ker{\mathop{\mathrm{Ker}}}
\def\Im{\mathop{\mathrm{Im}}}
\def\Coker{\mathop{\mathrm{Coker}}}
\def\Coim{\mathop{\mathrm{Coim}}}

% Boxtimes
%
\DeclareMathSymbol{\boxtimes}{\mathbin}{AMSa}{"02}

%
% Macros for moduli stacks/spaces
%
\def\QCohstack{\mathcal{QC}\!\mathit{oh}}
\def\Cohstack{\mathcal{C}\!\mathit{oh}}
\def\Spacesstack{\mathcal{S}\!\mathit{paces}}
\def\Quotfunctor{\mathrm{Quot}}
\def\Hilbfunctor{\mathrm{Hilb}}
\def\Curvesstack{\mathcal{C}\!\mathit{urves}}
\def\Polarizedstack{\mathcal{P}\!\mathit{olarized}}
\def\Complexesstack{\mathcal{C}\!\mathit{omplexes}}
% \Pic is the operator that assigns to X its picard group, usage \Pic(X)
% \Picardstack_{X/B} denotes the Picard stack of X over B
% \Picardfunctor_{X/B} denotes the Picard functor of X over B
\def\Pic{\mathop{\mathrm{Pic}}\nolimits}
\def\Picardstack{\mathcal{P}\!\mathit{ic}}
\def\Picardfunctor{\mathrm{Pic}}
\def\Deformationcategory{\mathcal{D}\!\mathit{ef}}
\usepackage{cleveref}
\crefname{originaldefinition}{original definition}{original definitions}
\Crefname{originaldefinition}{Original definition}{Original definitions}
\makeatletter
\newlabel{schemes-definition-reduced-induced-scheme}{{01J4}{source}{}{sourcepointer-1}{}}
\newlabel{schemes-definition-reduced-induced-scheme@cref}{{[originaldefinition][0][]01J4}{[1][0][]source}{}{}{}}
\makeatother
\crefname{originallemma}{original lemma}{original lemmas}
\Crefname{originallemma}{Original lemma}{Original lemmas}
\makeatletter
\newlabel{lemma-map-groupoids-on-perfect-field-image}{{04Q8}{source}{}{sourcepointer-2}{}}
\newlabel{lemma-map-groupoids-on-perfect-field-image@cref}{{[originallemma][0][]04Q8}{[1][0][]source}{}{}{}}
\makeatother
\crefname{originallemma}{original lemma}{original lemmas}
\Crefname{originallemma}{Original lemma}{Original lemmas}
\makeatletter
\newlabel{varieties-lemma-perfect-reduced}{{020I}{source}{}{sourcepointer-3}{}}
\newlabel{varieties-lemma-perfect-reduced@cref}{{[originallemma][0][]020I}{[1][0][]source}{}{}{}}
\makeatother
\crefname{originallemma}{original lemma}{original lemmas}
\Crefname{originallemma}{Original lemma}{Original lemmas}
\makeatletter
\newlabel{varieties-lemma-geometrically-reduced-any-base-change}{{035Z}{source}{}{sourcepointer-4}{}}
\newlabel{varieties-lemma-geometrically-reduced-any-base-change@cref}{{[originallemma][0][]035Z}{[1][0][]source}{}{}{}}
\makeatother
\crefname{originallemma}{original lemma}{original lemmas}
\Crefname{originallemma}{Original lemma}{Original lemmas}
\makeatletter
\newlabel{lemma-groupoid-on-field-image}{{04QB}{source}{}{sourcepointer-5}{}}
\newlabel{lemma-groupoid-on-field-image@cref}{{[originallemma][0][]04QB}{[1][0][]source}{}{}{}}
\makeatother
\title{The reduced closure of the image of j = (t, s) carries the restricted groupoid structure over a perfect field}
\author{The Stacks Project Authors}
\date{stacks-project:more-groupoids.tex:local-snapshot-20261012; GNU FDL 1.2 or later}
\begin{document}
\maketitle
\setcounter{subsection}{1}
\noindent\textit{Editorial note.} All delivered mathematical text is the author's own, copied from the pinned Stacks Project chapter file more-groupoids.tex; no statement or proof step was rewritten, reordered or removed. The counted span is the author environment and its own proof environment. References inside the span are handed to the assembler as navigation stubs that carry the canonical tag of the cited tag only; whether each of them is load-bearing is recorded on the stub itself.\par\medskip
\noindent\textit{Allowed standard prerequisites (editorial).}\begin{itemize}
\item Definition of the reduced closed subscheme on a closed subset (schemes chapter, Definition schemes-definition-reduced-induced-scheme, canonical tag 01J4; kind definition). cited at statement lines 1930-1932. Defines Z as the reduced closed subscheme on the closure of the image of j.
\item The map-groupoids-on-perfect-field-image lemma, of which the groupoid case is a special instance (more-groupoids chapter, Lemma lemma-map-groupoids-on-perfect-field-image, canonical tag 04Q8). cited at proof lines 1943-1945. One route to the result: the ambient groupoid makes the statement a special case.
\item A reduced scheme over a perfect field is geometrically reduced (varieties chapter, Lemma varieties-lemma-perfect-reduced, canonical tag 020I). cited at proof lines 1949-1952. Makes Z geometrically reduced over k via either projection, the entry point of the reducedness chain.
\item A geometrically reduced scheme stays reduced after any field-extension base change (varieties chapter, Lemma varieties-lemma-geometrically-reduced-any-base-change, canonical tag 035Z). cited at proof lines 1953-1955. Shows Z times\_\{pr\_1, U, pr\_0\} Z inside Z times Z is reduced, so the set theoretic containment of the field-image lemma is carried by reduced subschemes.
\item Set theoretic composition containment for the reduced closure of the image of j over a field (more-groupoids chapter, Lemma lemma-groupoid-on-field-image, canonical tag 04QB; collected on this lane as L22522). cited at proof lines 1956-1959. Provides pr\_02(Z times Z) inside Z set theoretically, which by reducedness upgrades to the induced morphism of the restricted fibre product into Z.
\end{itemize}\medskip
\ifdefined\hypertarget\hypertarget{sourcepointer-1}{}\fi
\noindent\textit{Source reference.} Original definition 01J4: sources/stacks-chow/tags.txt, source lines 2001--2001. Cited at statement lines 1930-1932: defines Z as the reduced closed subscheme on the closure of the image of j.\par\medskip
\ifdefined\hypertarget\hypertarget{sourcepointer-2}{}\fi
\noindent\textit{Source reference.} Original lemma 04Q8: sources/stacks-chow/tags.txt, source lines 6141--6141. Cited at proof lines 1943-1945: the ambient groupoid makes the statement a special case of this lemma.\par\medskip
\ifdefined\hypertarget\hypertarget{sourcepointer-3}{}\fi
\noindent\textit{Source reference.} Original lemma 020I: sources/stacks-chow/tags.txt, source lines 2635--2635. Cited at proof lines 1949-1952: Z is geometrically reduced over the perfect field k via either projection.\par\medskip
\ifdefined\hypertarget\hypertarget{sourcepointer-4}{}\fi
\noindent\textit{Source reference.} Original lemma 035Z: sources/stacks-chow/tags.txt, source lines 4146--4146. Cited at proof lines 1953-1955: Z times\_\{pr\_1, U, pr\_0\} Z inside Z times Z is reduced.\par\medskip
\ifdefined\hypertarget\hypertarget{sourcepointer-5}{}\fi
\noindent\textit{Source reference.} Original lemma 04QB: sources/stacks-chow/tags.txt, source lines 6144--6144. Cited at proof lines 1956-1959: the set theoretic containment pr\_02(Z times Z) inside Z that reducedness upgrades to the induced morphism.\par\medskip
\begin{lemma}
\label{lemma-groupoid-on-perfect-field-image}
Let $S$ be a scheme. Let $(U, R, s, t, c)$ be a groupoid scheme
over $S$. Assume $U$ is the spectrum of a perfect field.
Let $Z \subset U \times_S U$ be the reduced closed subscheme (see
Schemes, Definition \ref{schemes-definition-reduced-induced-scheme})
whose underlying topological space is the closure of the image of
$j = (t, s) : R \to U \times_S U$.
Then
$$
(U, Z, \text{pr}_0|_Z, \text{pr}_1|_Z,
\text{pr}_{02}|_{Z \times_{\text{pr}_1, U, \text{pr}_0} Z})
$$
is a groupoid scheme over $S$.
\end{lemma}

\begin{proof}
As $(U, U \times_S U, \text{pr}_1, \text{pr}_0, \text{pr}_{02})$
is a groupoid scheme over $S$ this is a special case of
Lemma \ref{lemma-map-groupoids-on-perfect-field-image}.
But we can also prove it directly as follows.

\medskip\noindent
We first explain why the statement makes sense. Since $U$ is the spectrum
of a perfect field $k$, the scheme $Z$ is geometrically reduced
over $k$ (via either projection), see
Varieties, Lemma \ref{varieties-lemma-perfect-reduced}.
Hence the scheme $Z \times_{\text{pr}_1, U, \text{pr}_0} Z \subset Z$
is reduced, see
Varieties, Lemma \ref{varieties-lemma-geometrically-reduced-any-base-change}.
Hence by
Lemma \ref{lemma-groupoid-on-field-image}
we see that $\text{pr}_{02}$ induces a morphism
$Z \times_{\text{pr}_1, U, \text{pr}_0} Z \to Z$.
Finally, it is clear that $\Delta_{U/S}$ factors through $Z$
and that the map
$\sigma : U \times_S U \to U \times_S U$, $(x, y) \mapsto (y, x)$
preserves $Z$. Since
$(U, U \times_S U, \text{pr}_0, \text{pr}_1, \text{pr}_{02},
\Delta_{U/S}, \sigma)$
satisfies the axioms of a groupoid, it follows that after restricting
to $Z$ they satisfy the axioms.
\end{proof}

\end{document}
